\section{Introduction}

Automatic summarization is one of the central problems in Natural
Language Processing (NLP) posing several challenges relating to {\em
  understanding} (i.e.,~identifying important content) and {\em
  generation} (i.e.,~aggregating and rewording the identified
content into a summary). Of the many summarization paradigms that have
been identified over the years (see \citeauthor{mani2001automatic},
\citeyear{mani2001automatic} and \citeauthor{Nenkova:McKeown:2011},
\citeyear{Nenkova:McKeown:2011} for a comprehensive overview),
single-document summarization has consistently attracted attention
\cite{jp-acl16,durrett-nyt-ext,nallapati-signll16,nallapati17,see-acl17,tanwan-acl17,narayan-arxiv17,Fan2017Controllable,paulus-socher-arxiv17,Pasunuru-multireward18,asli-multiagent18,narayan-sidenet18,narayan-rank18}.


\begin{figure}[t!]
  \center{\fontsize{10}{12}\selectfont % \footnotesize % \small % scriptsize % footnotesize
    \begin{tabular}{ p{7.2cm} }
      \thickhline
      
     % \vspace{0.01cm}
      
      \textbf{\textsc{Summary:}} \emph{\textcolor{orangered}{A man and
          a child have been killed} after a \textcolor{sepia}{light}
        \textcolor{forestgreen}{aircraft made an emergency landing} on
        \textcolor{midnightblue}{a beach in Portugal}.}  
       \\ \thickhline
      
      \textbf{\textsc{Document:}} Authorities said the incident took place on \textcolor{midnightblue}{Sao Joao beach in Caparica, south-west of Lisbon}. \\ 

      The National Maritime Authority said \textcolor{orangered}{a middle-aged man and a young girl died} after they were unable to avoid the plane. \\
      
      [{\em 6 sentences with 139 words are abbreviated from here.}] \\
      
      % The plane's only two occupants were unharmed, it added. \\

      % The Diario de Noticias newspaper quoted an eyewitness who said the plane had been flying at a low altitude over the beach, although he did not realise anything was wrong until other beachgoers began running. \\

      % One young witness told Reuters news agency: "I was near the water when I saw the plane. I called my parents, the plane fell on the sand and ran over two people, fatally hurting them and another was injured, I think, but I'm not sure, people were running away. \\

      % "The plane is still there, but the ambulances arrived quickly. I think maybe the fuel ran out because I find it weird that it landed on the beach." \\

      Other reports said the victims had been sunbathing when \textcolor{forestgreen}{the plane made its emergency landing}. \\

      [{\em Another 4 sentences with 67 words are abbreviated from here.}] \\

      % The Associated Press news agency said the girl who died had been with her parents, who were unhurt. The agency quoted witnesses from local television broadcasts. \\

      % Joao Quadros, who was on the beach, tweeted photos of the aftermath, saying the plane had passed by his son by a matter of metres. There had been no noise, he said. \\ 

      Video footage from the scene carried by local broadcasters showed \textcolor{sepia}{a small recreational plane} parked on the sand, apparently intact and surrounded by beachgoers and emergency workers. % One wing seemed to be misaligned in those photos.
      \\

      % The cause of the emergency landing remains unclear.
      % \\
        
      [{\em Last 2 sentences with 19 words are abbreviated.}] \\
      
      \thickhline
    \end{tabular}     
  }
  \caption{An abridged example from our extreme summarization dataset
    showing the document and its one-line summary. Document content
    present in the summary is color-coded.
    %The summary exemplifies across-the-document reasoning to identify important content and 
    %multiple levels of abstraction such as sentence fusion % ,
                                                           % synthesis
    %and paraphrasing. 
    % Title: Two killed in Portugal beach plane crash
    % url: http://www.bbc.co.uk/news/world-europe-40806376
  }\label{fig:bbcex-1}
\end{figure}

Neural approaches to NLP and their ability to learn continuous
features without recourse to pre-processing tools or linguistic
annotations have driven the development of large-scale document
summarization datasets
\cite{nytcorpus,hermann-nips15,newsroom-naacl18}.
% and their \shaycomment{ability of what? approach? neural networks? grammatically confusing}
% ability to learn continuous features
% without recourse to preprocessing tools or linguistic annotations has
% driven the development and use of large scale single document
% summarization datasets
% \cite{nytcorpus,hermann-nips15,newsroom-naacl18}.
However, these datasets often favor extractive models which create a
summary by identifying (and subsequently concatenating) the most
important sentences in a document
\cite{jp-acl16,nallapati17,narayan-rank18}.  Abstractive approaches,
despite being more faithful to the actual summarization task, either
lag behind extractive ones or are mostly extractive, exhibiting a
small degree of abstraction
\cite{see-acl17,tanwan-acl17,paulus-socher-arxiv17,Pasunuru-multireward18,asli-multiagent18}.


In this paper we introduce \emph{extreme summarization}, a new
single-document summarization task which is not amenable to extractive
strategies and requires an abstractive modeling approach. The idea is
to create a short, one-sentence news summary answering the question
``What is the article about?''. An example of a document and its
extreme summary are shown in Figure~\ref{fig:bbcex-1}. As can be seen,
the summary is very different from a headline whose aim is to
encourage readers to read the story; it draws on information
interspersed in various parts of the document (not only the beginning)
and displays multiple levels of abstraction including paraphrasing,
fusion, synthesis, and inference. We build a dataset for the proposed
task by harvesting online articles from the British Broadcasting
Corporation (BBC) that often include a first-sentence summary.
%answering ``What is the article about?''  Extreme summarization aims
%at generating the first sentence from the rest of the article.





%address the issue of this ``extractive bias''.  We
%introduce a new single document summarization task, ``extreme
%summarization,'' which does not favor extractive strategies and
%essentially requires abstractive summarization. 

%Figure~\ref{fig:bbcex-1} shows an example article
%from our dataset. Information to assemble the summary comes from
%different depths of the document demanding multiple levels of
%abstractions including paraphrasing, fusion, synthesis and
%inference. By the virtue of this task, it presents more challenges
%than extractive summarization while upholding the important
%characteristics of NLP and document summarization.


We further propose a novel deep learning model which we argue is
well-suited to the extreme summarization task. Unlike most existing
abstractive approaches
\cite{rush-acl15,chenIjcai-16,nallapati-signll16,see-acl17,tanwan-acl17,paulus-socher-arxiv17,Pasunuru-multireward18,asli-multiagent18}
which rely on an encoder-decoder architecture modeled by recurrent
neural networks (RNNs), we present a \emph{topic-conditioned} neural
model which is based entirely on convolutional neural networks
\cite{convseq2seq}. Convolution layers capture long-range
dependencies between words in the document more effectively compared to
RNNs, allowing to perform document-level inference, abstraction, and
paraphrasing.  Our convolutional encoder associates each word with a
topic vector capturing whether it is representative of the document's
content, while our convolutional decoder conditions each word
prediction on a document topic vector.


%and that are prone to mishandling long-range
%dependencies, we present a topic-conditioned convolutional
%sequence-to-sequence which allows document-level inference,
%abstraction and paraphrasing.  
%
%Our models adapts an architecture based
%entirely on convolutional neural networks (\textsc{ConvS2S},
%\citeauthor{convseq2seq}, \citeyear{convseq2seq}). We stack
%convolution layers to model long-range dependencies between words in
%the document effectively; they provide a shorter path to capture
%relationship between distant words in the document compared to RNNs.



% words at closer distances interact at lower layers while distant
% words interact at higher layers.
%In addition, we improve the convolutional
%encoder by associating each word with a topic vector capturing ``how
%topical is a word in the document?'' and the convolutional decoder by
%conditioning each word prediction on the document topic vector.


Experimental results show that when evaluated automatically (in terms
of ROUGE) our topic-aware convolutional model outperforms an oracle
extractive system and state-of-the-art RNN-based abstractive
systems. We also conduct two human evaluations in order to assess (a)
which type of summary participants prefer and (b) how much key
information from the document is preserved in the summary. Both
evaluations overwhelmingly show that human subjects find our summaries
more informative and complete. Our contributions in this work are
three-fold: a new single document summarization dataset that
encourages the development of abstractive systems; corroborated by
analysis and empirical results showing that extractive approaches are
not well-suited to the extreme summarization task; and a novel
topic-aware convolutional sequence-to-sequence model for abstractive
summarization.
%and an additional novel topic-conditioned model for summarization
%(also based on convolutional seq2seq modeling).


%%% Local Variables:
%%% mode: latex
%%% TeX-master: "extreme-sum"
%%% End:

